\section{Evaluate the observed margin, then ask a causal question}
\label{s6:ancestral-target}
Hidden common causes can create adjacencies absent from the observed portion
of a structural DAG. Consequently precision against that structural portion
is not a measure of recovering the hidden-variable method's target. A
finite-window construction makes the distinction exact. Observe price,
uptake and balance at dates 0 and 1. At each date an unobserved liquidity
state causes price and balance; price causes uptake, which causes balance.
Each variable and the liquidity state have a persistence edge to date 1.
Independent initial innovations define the beginning of the experiment.
There is no silently discarded infinite past.

For this finite DAG $D$ and observed set $O$, join observed nodes $i,j$ if
no subset of $O\setminus\{i,j\}$ d-separates them in $D$. This gives the
adjacency of its maximal ancestral margin. Put a tail at $i$ when $i$ is an
ancestor of $j$ in $D$, and an arrowhead otherwise. There are no selection
variables, so no undirected selection edges are required. An adjacency with
a tail may represent an inducing path, not a directly observed structural
coefficient. A bidirected edge records neither endpoint being an ancestor
of the other. The relation to latent projection, ancestral graphs and
window boundaries is developed in \citet{zhang2008,gerhardus2021}.

This example has six observed nodes and twelve marginal adjacencies. The
oracle checks every one of the 240 pair/conditioning-set queries twice,
using independent NetworkX and Tigramite d-separation implementations
\citep{runtime_networkx_dynamic}.
They agree on all queries. With the resulting population independence
oracle, LPCMCI recovers all twelve adjacencies and makes twelve correct
definite endpoint claims, leaving twelve circles. Those circles are not
errors against a fully oriented generating MAG: a PAG describes an
equivalence class, and unidentifiable orientations must remain unresolved.
The comparison measures correctness of definite ancestry claims and their
informativeness separately. It does not assert completeness of every
orientation under additional temporal background knowledge.

The finite-data experiment uses 2,500 independent two-period episodes per
replication and unrolls their six nodes into columns. LPCMCI is called with
maximum lag zero on these episode rows; this isolates its latent graph
operation without asserting stationarity across the finite window. The
earlier stationary time-series runs remain separate evidence. Treating this
episode construction as an infinite time-series oracle would be wrong:
longer unobserved history can alter its marginal adjacencies, as the
time-series ancestral-graph analysis explains \citep{gerhardus2021}.

With partial-correlation tests, fixed significance .01 and conditioning cap
four, forty repetitions give mean adjacency precision .9977 and recall
.8313, but mean accuracy of definite endpoint claims only .6605. Conflicting
marks are retained as conflicts, not silently converted to arrowheads or
tails. The strong adjacency precision therefore does not qualify a causal
ordering. In the exact population covariance, some genuinely adjacent pairs
have partial correlations as small as .00123 and .00869 under allowed
conditioning sets. Finite testing can remove such edges and then misorient
others. This supplies a concrete weak-faithfulness difficulty; the oracle
comparison shows that changing the target definition alone does not fix
statistical graph recovery. No favorable significance threshold is selected
after these results. Learned endpoints remain exploratory unless supported
by additional scientific information.

\subsection{An identifiable hidden-variable query}
Take the first-date price $A$, uptake $M$ and balance outcome $Y$ under the
substantive front-door graph $A\to M\to Y$, $A\leftrightarrow Y$.
Assume no direct price-to-outcome path, no unblocked back-door path from
$A$ to $M$, and all back-door paths from $M$ to $Y$ blocked by $A$.
These are stronger restrictions than an arbitrary LPCMCI output. Graphical
identification gives \citep{pearl2009,id2006}
\begin{equation}
 \E[Y\mid\doop(A=a)]
 =\sum_mP(m\mid a)\sum_{a'}\E[Y\mid m,a']P(a').
 \label{s6:frontdoor}
\end{equation}
The first factor identifies the action's effect on the mediator. Conditioning
on $A$ closes the mediator/outcome back door, allowing the second factor to
average the mediator intervention over the original action distribution.
The outer sum composes these interventions because the mediator intercepts
every directed path from $A$ to $Y$. An arbitrary adjustment regression on
$M$ would instead remove the mediated effect.

For an exact binary check, let $U$ be fair binary,
$P(A=1\mid U)=.2+.6U$, $P(M=1\mid A)=.15+.6A$ and
$P(Y=1\mid M,U)=.1+.3M+.4U$, with independent assignment noises.
The observed $2\times2\times2$ probability table is enumerated, not fitted.
Equation~\eqref{s6:frontdoor} recovers the structural contrast
$.6\times.3=.18$, while the unadjusted association is .42. A servicing or
deposit application must justify that all relevant pathways truly pass
through the measured uptake variable; a price's direct cash cost would
violate that exclusion if it were included in $Y$. Financial accounting can
add a known direct cost separately, as in the funding dossier, but cannot
wish away an unknown direct behavioral pathway.

\subsection{An unresolved query and its information requirement}
For an action/outcome pair with unrestricted hidden confounding, compare
two SCMs: $U\sim\operatorname{Bernoulli}(1/2)$ and $A=U$ in both,
with $Y=A$ in one and $Y=U$ in the other. Both produce
$P(A=Y=0)=P(A=Y=1)=1/2$. Their average action effects are 1 and 0.
Thus even an exactly known observed distribution and adjacency do not
identify this query. A larger sample or a preferred completion cannot
distinguish the two models.

Additional feasible randomization, observed action-relevant confounders
with support, a defensible instrument with its own target, or validated
front-door restrictions may resolve a suitably redesigned question. Without
such information the product returns the ambiguity and the requested
additional evidence. It does not turn circles into a recursive order.
Fresh estimation data address reuse of outcomes after graph selection;
they do not make the selected graph causally correct. The hidden-variable
query operation therefore records whether identification comes from an
asserted scientific graph, an invariant equivalence-class conclusion or
an unresolved discovery result.
